\section{狭义相对论基本原理与洛伦兹变换}
\subsection{狭义相对论基本原理}
成立前提：在彼此相对作匀速直线运动的\textbf{惯性参考系}中。

\begin{enumerate}
	\item \textbf{相对论原理：}物理定律在一切惯性参考系中都具有相同的数学表达形式。（所有惯性参考系对于描述物理现象都是等价的）
	\item \textbf{光速不变原理：}在彼此相对作匀速直线运动的任一惯性参考系中，所测得的光在真空中的传播速度都是相等的。
\end{enumerate}

\subsection{洛伦兹变换}
反映了两个相对作匀速直线运动的惯性参考系中的时空变换关系。

变换关系满足
$$\begin{cases}
	x^{'}=\frac{x-ut}{\sqrt{1-(\frac{u}{c})^2}}\\
	y^{'}=y\\
	z^{'}=z\\
	t^{'}=\frac{t-\frac{ux}{c^2}}{\sqrt{1-(\frac{u}{c})^2}}
\end{cases}$$
或
$$\begin{cases}
	x=\frac{x^{'}+ut}{\sqrt{1-(\frac{u}{c})^2}}\\
	y=y^{'}\\
	z=z^{'}\\
	t=\frac{t^{'}+\frac{ux}{c^2}}{\sqrt{1-(\frac{u}{c})^2}}
\end{cases}$$
\begin{enumerate}
	\item 洛伦兹变换与伽利略变换：当$u\ll c$时，即比值$\beta=\frac{u}{c}$很小时，洛伦兹变换变换就转化为伽利略变换。
	\item 相对论并没有把经典力学推翻，只是揭示了它的局限性。
	\item 相对论指出物体的速度不能超过真空中的光速。
\end{enumerate}
